\documentclass[../main.tex]{subfiles}

\begin{document}

\chapter{Exercises and Solutions}\label{app:exercises}

This appendix provides exercises for each chapter, ranging from computational practice to theoretical proofs. Selected solutions are provided at the end of each section.

%% ============================================================
\section{Chapter 1: Convex Optimization}\label{sec:ex-convex}

\subsection{Exercises}

\begin{simpleexercise}[Convexity Proofs]
Prove that the following functions are convex:
\begin{enumerate}[(a)]
    \item $f(x) = e^{ax}$ for any $a \in \mathbb{R}$
    \item $f(x) = x \log x$ for $x > 0$
    \item $f(x) = \|Ax - b\|_2^2$ for any matrix $A$ and vector $b$
    \item $f(X) = -\log \det(X)$ for positive definite $X$
\end{enumerate}
\end{simpleexercise}

\begin{simpleexercise}[Convex Sets]
Determine whether the following sets are convex. Prove your answer.
\begin{enumerate}[(a)]
    \item $\{x \in \mathbb{R}^n : \|x\|_2 \leq 1\}$ (unit ball)
    \item $\{x \in \mathbb{R}^n : \|x\|_0 \leq k\}$ (at most $k$ nonzero entries)
    \item $\{X \in \mathbb{R}^{n \times n} : X \succeq 0\}$ (positive semidefinite cone)
    \item $\{x \in \mathbb{R}^n : Ax = b, x \geq 0\}$
\end{enumerate}
\end{simpleexercise}

\begin{simpleexercise}[Operations Preserving Convexity]
Let $f$ and $g$ be convex functions. Prove:
\begin{enumerate}[(a)]
    \item $\alpha f + \beta g$ is convex for $\alpha, \beta \geq 0$
    \item $\max\{f, g\}$ is convex
    \item $f(Ax + b)$ is convex for any matrix $A$ and vector $b$
    \item Is $\min\{f, g\}$ convex? Provide a proof or counterexample.
\end{enumerate}
\end{simpleexercise}

\begin{simpleexercise}[Strong Convexity]
A function $f$ is $\mu$-strongly convex if $f(x) - \frac{\mu}{2}\|x\|^2$ is convex.
\begin{enumerate}[(a)]
    \item Prove that $f(x) = \frac{1}{2}x^\top Q x$ is strongly convex with parameter $\mu = \lambda_{\min}(Q)$
    \item Show that strong convexity implies a unique minimizer
    \item Prove that if $f$ is $\mu$-strongly convex and $g$ is convex, then $f + g$ is $\mu$-strongly convex
\end{enumerate}
\end{simpleexercise}

\begin{simpleexercise}[Lagrangian Duality]
Consider the optimization problem:
\[
\min_{x \in \mathbb{R}^n} \frac{1}{2}\|x\|^2 \quad \text{subject to } Ax = b
\]
\begin{enumerate}[(a)]
    \item Write the Lagrangian and derive the dual problem
    \item Solve the dual problem explicitly
    \item Recover the primal solution from the dual solution
    \item Verify strong duality holds
\end{enumerate}
\end{simpleexercise}

\begin{simpleexercise}[KKT Conditions]
For the problem $\min_{x} f(x)$ subject to $g_i(x) \leq 0$ and $h_j(x) = 0$:
\begin{enumerate}[(a)]
    \item State the KKT conditions
    \item Apply KKT to solve: $\min x_1^2 + x_2^2$ subject to $x_1 + x_2 = 1$
    \item Apply KKT to solve: $\min x_1 + x_2$ subject to $x_1^2 + x_2^2 \leq 1$
\end{enumerate}
\end{simpleexercise}

\subsection{Selected Solutions}

\begin{solution}[Exercise 1.1(c)]
We need to show $f(x) = \|Ax - b\|_2^2$ is convex.

\textbf{Method 1 (Second derivative):}
\[
f(x) = (Ax - b)^\top(Ax - b) = x^\top A^\top A x - 2b^\top Ax + b^\top b
\]
The Hessian is $\nabla^2 f(x) = 2A^\top A$, which is positive semidefinite since for any $v$:
\[
v^\top (A^\top A) v = \|Av\|_2^2 \geq 0
\]

\textbf{Method 2 (Composition):}
The squared norm $\|\cdot\|_2^2$ is convex (prove via definition). Affine functions $Ax - b$ preserve convexity under composition with convex increasing functions.
\end{solution}

\begin{solution}[Exercise 1.5]
\textbf{(a)} The Lagrangian is:
\[
L(x, \nu) = \frac{1}{2}\|x\|^2 + \nu^\top(Ax - b)
\]

\textbf{(b)} Setting $\nabla_x L = x + A^\top\nu = 0$ gives $x = -A^\top\nu$. Substituting into the Lagrangian:
\[
g(\nu) = \frac{1}{2}\|A^\top\nu\|^2 - \nu^\top A A^\top\nu - \nu^\top b = -\frac{1}{2}\nu^\top AA^\top\nu - \nu^\top b
\]
The dual problem is $\max_\nu g(\nu) = \max_\nu -\frac{1}{2}\nu^\top AA^\top\nu - \nu^\top b$.

Setting $\nabla_\nu g = 0$: $-AA^\top\nu - b = 0$, so $\nu^* = -(AA^\top)^{-1}b$.

\textbf{(c)} The primal solution: $x^* = -A^\top\nu^* = A^\top(AA^\top)^{-1}b$.

\textbf{(d)} This is the minimum norm solution to $Ax = b$. Strong duality holds because the problem is convex with linear constraints (Slater's condition is satisfied).
\end{solution}

%% ============================================================
\section{Chapter 2: Stochastic Optimization}\label{sec:ex-stochastic}

\subsection{Exercises}

\begin{simpleexercise}[SGD Variance]
Consider $f(x) = \frac{1}{n}\sum_{i=1}^n f_i(x)$ where $f_i(x) = \frac{1}{2}(a_i^\top x - b_i)^2$.
\begin{enumerate}[(a)]
    \item Compute $\nabla f_i(x)$ and $\nabla f(x)$
    \item Derive an expression for $\mathbb{E}[\|\nabla f_i(x) - \nabla f(x)\|^2]$
    \item Show how this variance depends on the data heterogeneity
    \item Under what conditions is the variance zero?
\end{enumerate}
\end{simpleexercise}

\begin{simpleexercise}[Mini-batch SGD]
Let $g_B = \frac{1}{|B|}\sum_{i \in B} \nabla f_i(x)$ be a mini-batch gradient.
\begin{enumerate}[(a)]
    \item Prove $\mathbb{E}[g_B] = \nabla f(x)$ (unbiasedness)
    \item Show that $\text{Var}(g_B) = \frac{\sigma^2}{|B|}$ where $\sigma^2$ is the single-sample variance
    \item Derive the optimal batch size given a computational budget
    \item Compare wall-clock time for batch sizes $B = 1, 32, 256$ assuming linear parallelization
\end{enumerate}
\end{simpleexercise}

\begin{simpleexercise}[Importance Sampling]
Consider weighted sampling where $p_i \propto \|\nabla f_i(x)\|$.
\begin{enumerate}[(a)]
    \item Derive the importance-weighted gradient estimator
    \item Prove it is unbiased
    \item Show it has lower variance than uniform sampling
    \item Discuss practical challenges of this approach
\end{enumerate}
\end{simpleexercise}

\begin{simpleexercise}[Noise and Generalization]
\begin{enumerate}[(a)]
    \item Explain why SGD noise can help escape saddle points
    \item How does batch size affect the implicit regularization of SGD?
    \item Design an experiment to test whether larger batch sizes lead to sharper minima
\end{enumerate}
\end{simpleexercise}

\subsection{Selected Solutions}

\begin{solution}[Exercise 2.2(a,b)]
\textbf{(a)} By linearity of expectation:
\[
\mathbb{E}[g_B] = \mathbb{E}\Big[\frac{1}{|B|}\sum_{i \in B} \nabla f_i(x)\Big] = \frac{1}{|B|}\sum_{i \in B} \mathbb{E}[\nabla f_i(x)] = \nabla f(x)
\]
since each $i$ is drawn uniformly and $\mathbb{E}[\nabla f_i(x)] = \nabla f(x)$.

\textbf{(b)} For independent samples:
\begin{align}
\text{Var}(g_B) &= \text{Var}\Big(\frac{1}{|B|}\sum_{i \in B} \nabla f_i(x)\Big) \\
&= \frac{1}{|B|^2} \sum_{i \in B} \text{Var}(\nabla f_i(x)) \\
&= \frac{1}{|B|^2} \cdot |B| \cdot \sigma^2 = \frac{\sigma^2}{|B|}
\end{align}
This shows variance decreases linearly with batch size.
\end{solution}

%% ============================================================
\section{Chapter 3: Gradient Descent}\label{sec:ex-gradient-descent}

\subsection{Exercises}

\begin{simpleexercise}[Learning Rate Analysis]
For $f(x) = \frac{1}{2}x^\top Qx$ with $Q = \text{diag}(\lambda_1, \ldots, \lambda_n)$:
\begin{enumerate}[(a)]
    \item Write the gradient descent update for each coordinate
    \item Show the update is $x_i^{(k+1)} = (1 - \eta\lambda_i)x_i^{(k)}$
    \item Find the range of $\eta$ for convergence
    \item Find the optimal $\eta$ minimizing worst-case convergence
    \item How does the condition number $\kappa = \lambda_{\max}/\lambda_{\min}$ affect convergence?
\end{enumerate}
\end{simpleexercise}

\begin{simpleexercise}[Momentum]
For heavy ball momentum: $v_{t+1} = \beta v_t + \nabla f(x_t)$, $x_{t+1} = x_t - \eta v_t$
\begin{enumerate}[(a)]
    \item Rewrite as a second-order recurrence in $x_t$
    \item For $f(x) = \frac{1}{2}\lambda x^2$, find the characteristic equation
    \item Determine optimal $\eta$ and $\beta$ for fastest convergence
    \item Compare the rate to vanilla gradient descent
\end{enumerate}
\end{simpleexercise}

\begin{simpleexercise}[Nesterov Momentum]
Implement Nesterov accelerated gradient and compare to heavy ball momentum on:
\begin{enumerate}[(a)]
    \item A quadratic: $f(x) = \frac{1}{2}x^\top Q x$ with $\kappa = 100$
    \item Logistic regression on a synthetic dataset
    \item A neural network on MNIST (first 1000 samples)
\end{enumerate}
Plot loss curves and compare convergence behavior.
\end{simpleexercise}

\begin{simpleexercise}[Learning Rate Schedules]
Implement and compare:
\begin{enumerate}[(a)]
    \item Constant learning rate
    \item Step decay (halve every $k$ epochs)
    \item Exponential decay: $\eta_t = \eta_0 \cdot \gamma^t$
    \item Cosine annealing: $\eta_t = \eta_{\min} + \frac{1}{2}(\eta_0 - \eta_{\min})(1 + \cos(\pi t/T))$
    \item Warmup + decay
\end{enumerate}
Test on a 2-layer MLP trained on CIFAR-10.
\end{simpleexercise}

\subsection{Selected Solutions}

\begin{solution}[Exercise 3.1]
\textbf{(a,b)} The gradient is $\nabla f(x) = Qx$, so $(\nabla f(x))_i = \lambda_i x_i$. The update:
\[
x_i^{(k+1)} = x_i^{(k)} - \eta \lambda_i x_i^{(k)} = (1 - \eta\lambda_i)x_i^{(k)}
\]

\textbf{(c)} For convergence, we need $|1 - \eta\lambda_i| < 1$ for all $i$. This requires:
\[
-1 < 1 - \eta\lambda_i < 1 \implies 0 < \eta < \frac{2}{\lambda_i}
\]
For all eigenvalues: $0 < \eta < 2/\lambda_{\max}$.

\textbf{(d)} The convergence rate is $\max_i |1 - \eta\lambda_i|$. To minimize the worst case:
\[
\eta^* = \frac{2}{\lambda_{\max} + \lambda_{\min}}
\]
giving rate $\rho = \frac{\kappa - 1}{\kappa + 1}$.

\textbf{(e)} For large $\kappa$: $\rho \approx 1 - 2/\kappa$. The number of iterations to reduce error by factor $\epsilon$ is $O(\kappa \log(1/\epsilon))$.
\end{solution}

%% ============================================================
\section{Chapter 4: Advanced Optimizers}\label{sec:ex-advanced}

\subsection{Exercises}

\begin{simpleexercise}[Implement Adam]
Write a complete implementation of Adam from scratch.
\begin{enumerate}[(a)]
    \item Implement the basic Adam algorithm
    \item Add weight decay (AdamW variant)
    \item Add AMSGrad modification
    \item Test on logistic regression and compare to SGD with momentum
\end{enumerate}
\end{simpleexercise}

\begin{simpleexercise}[Optimizer Comparison]
Compare SGD, SGD+Momentum, AdaGrad, RMSprop, and Adam on:
\begin{enumerate}[(a)]
    \item Rosenbrock function (2D)
    \item A sparse linear regression problem
    \item ResNet-18 on CIFAR-10 (10 epochs)
\end{enumerate}
Create plots showing loss curves, gradient norm evolution, and final accuracy.
\end{simpleexercise}

\begin{simpleexercise}[AdaGrad Analysis]
For AdaGrad: $x_{t+1} = x_t - \frac{\eta}{\sqrt{G_t + \epsilon}} \odot g_t$ where $G_t = \sum_{s=1}^t g_s \odot g_s$:
\begin{enumerate}[(a)]
    \item Show that effective learning rates decrease over time
    \item Prove this is beneficial for sparse problems
    \item What is the main limitation for non-convex optimization?
    \item How does RMSprop address this limitation?
\end{enumerate}
\end{simpleexercise}

\begin{simpleexercise}[Learning Rate Warmup]
\begin{enumerate}[(a)]
    \item Explain why warmup is necessary for Adam with large learning rates
    \item Implement linear warmup for the first $k$ steps
    \item Compare training with and without warmup on a transformer model
\end{enumerate}
\end{simpleexercise}

\subsection{Selected Solutions}

\begin{solution}[Exercise 4.1(a)]
\begin{verbatim}
def adam(params, grads, m, v, t, lr=0.001, beta1=0.9, beta2=0.999, eps=1e-8):
    """
    Adam optimizer update.

    Args:
        params: list of parameters
        grads: list of gradients
        m: first moment estimates
        v: second moment estimates
        t: timestep (starting from 1)
        lr: learning rate
        beta1, beta2: decay rates
        eps: numerical stability

    Returns:
        updated params, m, v
    """
    for i in range(len(params)):
        # Update biased first moment estimate
        m[i] = beta1 * m[i] + (1 - beta1) * grads[i]

        # Update biased second moment estimate
        v[i] = beta2 * v[i] + (1 - beta2) * (grads[i] ** 2)

        # Bias correction
        m_hat = m[i] / (1 - beta1 ** t)
        v_hat = v[i] / (1 - beta2 ** t)

        # Update parameters
        params[i] = params[i] - lr * m_hat / (np.sqrt(v_hat) + eps)

    return params, m, v
\end{verbatim}
\end{solution}

%% ============================================================
\section{Chapter 5: Generalization and Regularization}\label{sec:ex-regularization}

\subsection{Exercises}

\begin{simpleexercise}[L1 vs L2 Regularization]
Consider $\min_w \|Xw - y\|^2 + \lambda R(w)$ with $R(w) = \|w\|_1$ or $R(w) = \|w\|_2^2$.
\begin{enumerate}[(a)]
    \item Derive the gradient (or subgradient) for each case
    \item Solve analytically for $X = I$ (identity matrix)
    \item Explain geometrically why L1 promotes sparsity
    \item Implement and compare on a sparse signal recovery problem
\end{enumerate}
\end{simpleexercise}

\begin{simpleexercise}[Elastic Net]
Elastic net combines L1 and L2: $R(w) = \alpha\|w\|_1 + (1-\alpha)\|w\|_2^2$.
\begin{enumerate}[(a)]
    \item Derive the proximal operator for elastic net
    \item When is elastic net preferred over pure L1 or L2?
    \item Implement coordinate descent for elastic net regression
\end{enumerate}
\end{simpleexercise}

\begin{simpleexercise}[Dropout]
\begin{enumerate}[(a)]
    \item Derive the expected output of a dropout layer with rate $p$
    \item Explain why we scale by $1/(1-p)$ during training (inverted dropout)
    \item Show that dropout is equivalent to a form of weight regularization
    \item Compare dropout to L2 regularization on a 3-layer MLP
\end{enumerate}
\end{simpleexercise}

\begin{simpleexercise}[Early Stopping]
\begin{enumerate}[(a)]
    \item Prove that early stopping is equivalent to L2 regularization for linear regression with gradient descent
    \item Implement early stopping with a patience parameter
    \item How do you choose the validation set size?
    \item Design an experiment comparing early stopping to explicit regularization
\end{enumerate}
\end{simpleexercise}

\begin{simpleexercise}[Weight Decay vs L2]
\begin{enumerate}[(a)]
    \item For SGD, show that weight decay and L2 regularization are equivalent
    \item For Adam, show they are NOT equivalent
    \item Explain why AdamW (decoupled weight decay) often works better than Adam + L2
\end{enumerate}
\end{simpleexercise}

\subsection{Selected Solutions}

\begin{solution}[Exercise 5.1(b,c)]
\textbf{(b)} For $X = I$, the problem becomes:
\[
\min_w \frac{1}{2}\|w - y\|^2 + \lambda R(w)
\]
This has a closed-form solution:
\begin{itemize}
    \item \textbf{L2:} $w^* = \frac{y}{1 + 2\lambda}$ (uniform shrinkage)
    \item \textbf{L1:} $w_i^* = \text{sign}(y_i)(|y_i| - \lambda)_+$ (soft thresholding)
\end{itemize}

\textbf{(c)} The L1 ball $\{w : \|w\|_1 \leq t\}$ has corners on the axes. The level sets of the quadratic loss are ellipsoids. For generic data, these intersect at a corner, giving a sparse solution. The L2 ball is smooth, so intersections typically occur away from axes.
\end{solution}

%% ============================================================
\section{Chapter 6: Loss Functions}\label{sec:ex-loss}

\subsection{Exercises}

\begin{simpleexercise}[Cross-Entropy Loss]
\begin{enumerate}[(a)]
    \item Derive cross-entropy loss from maximum likelihood estimation
    \item Show that minimizing cross-entropy is equivalent to minimizing KL divergence
    \item Compute the gradient of cross-entropy with respect to logits
    \item Why do we use log-softmax for numerical stability?
\end{enumerate}
\end{simpleexercise}

\begin{simpleexercise}[Loss Function Design]
Design appropriate loss functions for:
\begin{enumerate}[(a)]
    \item Predicting house prices (regression with outliers)
    \item Multi-label classification (each sample can have multiple labels)
    \item Ordinal regression (predicting ratings 1-5)
    \item Learning to rank (relative ordering matters)
\end{enumerate}
\end{simpleexercise}

\begin{simpleexercise}[Focal Loss]
Focal loss modifies cross-entropy: $FL(p_t) = -(1 - p_t)^\gamma \log(p_t)$.
\begin{enumerate}[(a)]
    \item Show that $\gamma = 0$ recovers cross-entropy
    \item Plot focal loss for $\gamma = 0, 1, 2, 5$
    \item Explain why focal loss helps with class imbalance
    \item Implement and test on an imbalanced dataset
\end{enumerate}
\end{simpleexercise}

\begin{simpleexercise}[Huber Loss]
The Huber loss is:
\[
L_\delta(y, \hat{y}) = \begin{cases}
\frac{1}{2}(y - \hat{y})^2 & |y - \hat{y}| \leq \delta \\
\delta(|y - \hat{y}| - \frac{\delta}{2}) & \text{otherwise}
\end{cases}
\]
\begin{enumerate}[(a)]
    \item Verify that Huber loss is continuous and differentiable
    \item Show it behaves like MSE near zero and MAE for large errors
    \item Compute the gradient
    \item Compare robustness to outliers vs MSE
\end{enumerate}
\end{simpleexercise}

\subsection{Selected Solutions}

\begin{solution}[Exercise 6.1(a,b)]
\textbf{(a)} For classification with true label $y$ and predicted probabilities $p$:
\[
\mathcal{L} = -\log P(y | x) = -\log p_y = -\sum_{k} y_k \log p_k
\]
where $y$ is one-hot encoded. This is cross-entropy.

\textbf{(b)} KL divergence from true distribution $q$ to predicted $p$:
\[
D_{KL}(q \| p) = \sum_k q_k \log \frac{q_k}{p_k} = -H(q) - \sum_k q_k \log p_k
\]
Since $H(q)$ (entropy of true distribution) is constant with respect to $p$, minimizing KL divergence is equivalent to minimizing cross-entropy.
\end{solution}

%% ============================================================
\section{Chapter 7: Convergence Analysis}\label{sec:ex-convergence}

\subsection{Exercises}

\begin{simpleexercise}[Convergence Rate Analysis]
For $f(x) = \frac{1}{2}x^\top Qx$ with eigenvalues $\mu \leq \lambda_i \leq L$:
\begin{enumerate}[(a)]
    \item Prove the convergence rate is $\rho = \frac{L - \mu}{L + \mu}$ for optimal step size
    \item How many iterations are needed to reduce error by factor $10^{-6}$?
    \item Compare to accelerated gradient descent rate $\rho_{acc} = \frac{\sqrt{L} - \sqrt{\mu}}{\sqrt{L} + \sqrt{\mu}}$
\end{enumerate}
\end{simpleexercise}

\begin{simpleexercise}[Condition Number]
\begin{enumerate}[(a)]
    \item Compute the condition number of $Q = \begin{pmatrix} 100 & 0 \\ 0 & 1 \end{pmatrix}$
    \item Plot gradient descent trajectories for $f(x) = \frac{1}{2}x^\top Q x$ with various condition numbers
    \item How does preconditioning improve the condition number?
    \item Design a preconditioner for sparse linear systems
\end{enumerate}
\end{simpleexercise}

\begin{simpleexercise}[Non-Convex Analysis]
For non-convex $f$ with $L$-Lipschitz gradient:
\begin{enumerate}[(a)]
    \item Prove that GD with $\eta = 1/L$ satisfies: $\min_{t \leq T} \|\nabla f(x_t)\|^2 \leq \frac{2(f(x_0) - f^*)}{T/L}$
    \item What does this say about convergence to stationary points?
    \item How do saddle points affect this analysis?
    \item Discuss the role of noise in escaping saddle points
\end{enumerate}
\end{simpleexercise}

\begin{simpleexercise}[Variance Reduction]
Implement and compare:
\begin{enumerate}[(a)]
    \item Vanilla SGD
    \item SVRG (Stochastic Variance Reduced Gradient)
    \item SAGA
\end{enumerate}
Test on logistic regression with $n = 10000$ samples. Plot variance of gradient estimates over time.
\end{simpleexercise}

\begin{simpleexercise}[Practical Convergence Diagnosis]
For a neural network training run:
\begin{enumerate}[(a)]
    \item Design metrics to monitor convergence (beyond just loss)
    \item Implement gradient norm tracking
    \item Detect signs of: divergence, saddle points, local minima
    \item Create visualizations for diagnosing optimization issues
\end{enumerate}
\end{simpleexercise}

\subsection{Selected Solutions}

\begin{solution}[Exercise 7.1(a,b)]
\textbf{(a)} With optimal step size $\eta^* = \frac{2}{L + \mu}$:
\[
x^{(k+1)} = x^{(k)} - \eta^* Qx^{(k)} = (I - \eta^* Q)x^{(k)}
\]
The convergence rate is:
\[
\rho = \|I - \eta^* Q\| = \max_i |1 - \eta^* \lambda_i| = \frac{L - \mu}{L + \mu} = \frac{\kappa - 1}{\kappa + 1}
\]

\textbf{(b)} For error reduction of $10^{-6}$:
\[
\rho^k \leq 10^{-6} \implies k \geq \frac{6 \ln 10}{\ln(1/\rho)} \approx \frac{6 \ln 10}{\frac{2}{\kappa}} = 3\kappa \ln 10 \approx 6.9\kappa
\]
For $\kappa = 100$, we need approximately $k \approx 690$ iterations.
\end{solution}

\end{document}
